The previous chapter demonstrated a feedback loop that works: CARP reconstructs a global view of
an evolving key distribution and adapts range partitioning online.
This chapter examines what it costs when the corresponding loop is high-latency, using our
experience developing telemetry-driven placement policies for block-structured adaptive mesh
refinement (AMR) codes as a case study.
AMR accelerates simulations by refining the mesh only where needed, but that adaptivity produces
heterogeneous per-block work and shifting communication
patterns~\cite{Dubey2014:Survey, Burst2011:p4est}.
At scale, frequent synchronization amplifies this variability into
stragglers, so placement and load balancing become first-order performance
concerns~\cite{Dubey2014:Survey, Burst2008:ALPS}.
Careful end-to-end characterization of these effects in production AMR codes is sparse.
The best available evidence from exported traces of large MPI
applications~\cite{Klenk2017:OverviewMPICharacteristics} reports substantial waiting time at scale
and emphasizes that standard trace data is often insufficient for cross-stack diagnosis, leaving
algorithmic and infrastructural causes entangled.
We therefore study real AMR codes on our own research cluster with full-stack access, collect our
own telemetry, and evaluate changes empirically.

Our goal was to replace static placement heuristics with telemetry-driven workload models, but
initial experiments were dominated by effects that had nothing to do with placement.
Fail-slow hardware~\cite{Gunaw2018:FailSlow}, fabric pathologies, MPI runtime
interactions, and cross-layer tuning issues produced straggler signatures that could not be
distinguished from genuine placement effects until they were identified and removed.
Localizing these behaviors required fine-grained telemetry that, in practice, was only accessible
through traces.
Standard tracing tools~\cite{Shend2006:TAU, Knupf2012:Score-P} collected this data as per-rank
logs, but the resulting volume was hard to control across thousands of timesteps, and the fixed
schemas did not capture the custom fields we needed to interpret AMR behavior.
We therefore iterated through a series of bespoke instrumentation and analysis pipelines,
repeatedly adjusting what to emit so that post-hoc analysis remained tractable.
The result was a slow, manual diagnostic cycle in which time was dominated by collecting, moving,
and reshaping telemetry, even when the eventual remediation was simple (for example, pruning bad
nodes, correcting configuration, or fixing an obvious tuning mistake).

With a reliable measurement baseline established, we could finally isolate placement effects.
Validated telemetry confirmed that synchronization dominated runtime (35--50\% of total time,
increasing with scale), driven by compute variability across blocks.
We found that compute load balance and communication locality are competing objectives: improving
one degrades the other.
We developed CPLX (\S\ref{sec:amr:design}), a hybrid placement policy that exposes this tradeoff
through a single tunable parameter $X$.
At $X{=}0$, CPLX prioritizes locality; at $X{=}100$, it reduces to pure load balancing.
On a real AMR code across 512--4096 CPU ranks, CPLX improves runtime by up to 21.6\% over tuned
baselines (\S\ref{sec:amr:eval}), with the best setting of $X$ varying with the code, problem,
scale, and cluster.

The entire exercise---from diagnosing thermal throttling to evaluating CPLX tradeoffs---was one
continuous cycle of hypothesis-driven telemetry work.
There was no clean boundary between denoising and placement.
Every intervention required re-collecting data, re-analyzing it, and re-validating that the change
produced the effect we believed it did.
The cost was establishing measurement fidelity, attributing variability to root causes across the
stack, and iterating toward a reliable baseline from which placement effects could be studied.
Developing and evaluating CPLX required the same pattern, including custom instrumentation for
per-block costs, post-hoc analysis, and repeated checks that apparent improvements reflected
placement effects rather than residual noise in the measurement and execution environment.
Every finding was specific to one code on one cluster, and the conditions that made the process
expensive---including cross-stack confounders, context-dependent tradeoffs, and reshaping telemetry
for each new question---are structural properties of large-scale parallel systems.
Because the dominant effects and the best tradeoffs depend on the code, problem, scale, and
environment, translating a one-off win into a repeatable practice requires repeating much of the
same process.

This chapter traces that process end to end: the confounders we encountered and how we diagnosed
them (\S\ref{sec:amr:chal}), the evolution of our telemetry and analysis workflow
(\S\ref{sec:amr:tele}), and the design and evaluation of CPLX
(\S\ref{sec:amr:design}, \S\ref{sec:amr:eval}).
It concludes with empirical guidance on what effective, iterative performance-diagnosis tooling
would need to support, based on what succeeded and what failed in our experience
(\S\ref{sec:amr:lessons}).

% Lessons for ORCA: and I think pt 2 should be:

% everything is context-dependent and requires reactions parsing CSVs in pandas was always painful.
% subsequent loading was nice if we wrote to feather (why were we writing csvs? we were just
% following established narratives, we had no idea what tools should we be using) nothing is
% constant. you can't assume canned views. every problem requried a different view. you can't assume
% a data model. tracing tools forced standard columns but we wanted block IDs to know which kernel
% belonged to which block. your telemetry data models need to be versatile. SQL works. we tried
% ETL-ing to clickhouse then triggering sql queries to get dataframes that we analyzed in pandas.
% that's the building block for the real thing.
% ebpf showed us the view-driven paradigm. when we were down in the weeds and everything was
% failing, ebpf was like a savior. instead of giving us telemetry, it asked us: all telemetry is
% available on top. tell me what do you want to see. this is the thing that needs to exist at
% cluster scale.
